\chapter*{Words To Know}
\markboth{Words To Know}{Words To Know}
\phantomsection
\addcontentsline{toc}{chapter}{Words To Know}

\begin{description}
\item[Animation] A trick that shows different pictures quickly so something
appears to move.

\item[BASIC] The language used in this book. NovaBASIC uses numbered lines and
commands such as \texttt{PRINT}, \texttt{GOTO}, and \texttt{FOR}.

\item[Bit] A tiny piece of information that is either on or off.

\item[Byte] Eight bits. Many Nova hardware values fit in one byte.

\item[Command] A word that tells BASIC to do something, such as \texttt{PRINT}
or \texttt{LINE}.

\item[Framework] Shared program code used by more than one game. In this book,
the common framework starts at line 60000.

\item[Graphics screen] The part of the display used for drawing pixels, lines,
rectangles, circles, sprites, roads, maps, and game boards.

\item[Line number] The number at the start of a BASIC line. NovaBASIC uses line
numbers to keep the program in order and to jump with \texttt{GOTO} or
\texttt{GOSUB}.

\item[Loop] A part of a program that repeats. Games use loops to keep reading
keys, moving objects, and redrawing the screen.

\item[Sprite] A small hardware-controlled picture that can move around the
screen without redrawing the whole background.

\item[String] Text stored in a variable. String variable names end with
\texttt{\$}, such as \texttt{NM\$}.

\item[Variable] A named place to store a value. In this book, variable names
are usually one or two characters long.

\item[XRAM] Nova expansion memory. Bigger programs can use it to store tables,
levels, pictures, and other data outside normal BASIC memory.
\end{description}
